\lyseis{Β3.1}{Στοιχεία και είδη τριγώνων}
\sol{1} α) ισόπλευρο β) ισοσκελές γ) σκαληνό.
\sol{2} $6{,}5+7+8{,}5=22$ cm.
\sol{3} Ελεύθερη κατασκευή. Στο ορθογώνιο δύο ύψη είναι οι κάθετες πλευρές, στο αμβλυγώνιο δύο ύψη βρίσκονται έξω από το τρίγωνο.
\sol{4} $\mathrm{ΒΜ}=4{,}4:2=2{,}2$ cm.
\sol{5} Οι τρεις διάμεσοι περνούν από το ίδιο σημείο.
\sol{6} Ελεύθερη κατασκευή. Οι δύο γωνίες είναι ίσες.
\sol{7} α) Σωστό. β) Λάθος. γ) Σωστό. δ) Σωστό. ε) Λάθος.
\sol{8} Ισοσκελές, περίμετρος $30+30+40=100$ cm.
\sol{9} $3\cdot60=180$ cm.
\sol{10} $24:3=8$ μ.
\sol{11} $(50-20):2=15$ cm.
\sol{12} Με $3,4,9$ όχι (γιατί $3+4=7<9$). Με $3,4,5$ ναι ($3+4=7>5$ κ.λπ.).
\sol{13} Οξυγώνιο, ορθογώνιο ή αμβλυγώνιο.
\sol{14} Όχι: ένα τρίγωνο με ορθή γωνία έχει τις άλλες δύο οξείες, άρα πλευρές όχι όλες ίσες. Το ισόπλευρο έχει όλες τις γωνίες $60^\circ$.
\sol{15} Δύο (στο αμβλυγώνιο).
\sol{16} Δεν υπάρχει ισόπλευρο ορθογώνιο ή αμβλυγώνιο. Υπάρχουν: ορθογώνιο σκαληνό, ορθογώνιο ισοσκελές, αμβλυγώνιο σκαληνό, αμβλυγώνιο ισοσκελές, οξυγώνιο σκαληνό, οξυγώνιο ισοσκελές, οξυγώνιο ισόπλευρο.
\sol{17} Ελεύθερη απάντηση.
\sol{18} Ελεύθερη κατασκευή.
